\documentclass[10pt,a4paper]{amsart}
\usepackage[margin=19mm]{geometry}
\usepackage{amsmath,amssymb,mathtools}
\usepackage[scr=rsfs,cal=euler]{mathalfa}
\usepackage{xfrac}
\usepackage[unicode,hidelinks,bookmarks=false]{hyperref}
\newcommand{\ie}{i.e.\ }
\newcommand{\cf}{cf.\ }
\newcommand{\resp}{respectively\ }
\newcommand{\Subsection}{\subsection}
\providecommand{\nobreakdash}{\nobreak\hbox{-}}
\setlength{\emergencystretch}{2em}
\multlinegap0pt
\usepackage{bm}
\usepackage{bbm}
\usepackage{dsfont}
\usepackage{xspace}
\usepackage{cancel}
\usepackage{mathtools}
\usepackage{amsthm}
\makeatletter
\@ifundefined{theorem}{\newtheorem{theorem}{Theorem}}{}
\@ifundefined{proposition}{\newtheorem{proposition}[theorem]{Proposition}}{}
\makeatother
\newcommand{\et}[1]{\mathrm e\!\left(#1\right)}
\title[Almost everywhere divergence of double Fourier series along ]{Almost everywhere divergence of double Fourier series along shrinking conical regions}
\author{Ushangi Goginava}
\date{Source: arXiv:2606.04119v1 (CC BY 4.0)}
\begin{document}
\maketitle

\noindent\emph{Source and licence.} Almost everywhere divergence of double Fourier series along shrinking conical regions. Version arXiv:2606.04119v1 (CC BY 4.0), distributed under the Creative Commons Attribution 4.0 International licence, as recorded in the arXiv licence metadata for this exact version and shown on the abstract page https://arxiv.org/abs/2606.04119v1. Full rights notice: licenses/arxiv-2606.04119-CC-BY-4.0.txt.\par\smallskip

\noindent\textit{Editorial scope.} Counted unit: the Proposition whose source label is \texttt{prop:block} in the article (source0.tex lines 226--248, source label \texttt{prop:block}), delivered together with the article's own complete original proof of it (source0.tex lines 250--451, one \texttt{proof} environment of 202 lines). No mathematics is written, repaired or re-derived: statement and proof are the article's own text line for line. Required context retained uncounted: eq:S-R-relation (lines 108--110). Every definition, display and label of those lines is retained unchanged. Omitted from this package and not redistributed: 1--107, 111--225, 249--249, 452--1087. Nothing inside a retained excerpt is omitted or altered: every retained body is delivered exactly as the article's lines are written, so no omission receipt of any kind accompanies this package. Citations: the retained excerpts carry no citation command at all, so no bibliography entry of the article is transcribed and no reference is redistributed. Reference adaptation: none was needed and none was made; every reference inside the delivered unit resolves inside it, so no unresolved reference or citation is produced. Printed slips kept literal: no printed word, symbol, hypothesis or number of the counted statement or of its proof was altered. Layout and class: the article's own class is \texttt{amsart} with packages including amsfonts, geometry, amsmath, amssymb, amsthm, mathtools, hyperref; the wrapper is the lane's pinned \texttt{amsart} editorial wrapper at 10pt on A4, which restates the theorem-like environments the retained text needs and adds the article's own macro definitions that the retained text needs (\texttt{\textbackslash et}), with extra wrapper packages (bm, bbm, dsfont, xspace, cancel, mathtools). The delivered numbering is the unit's own. Assets: no retained span contains a figure environment, a graphics-inclusion command, a PDF-page inclusion, a table or a TikZ picture; no image file is copied into this package, the package declares no asset dependency and no image file is redistributed with it. Licence: version arXiv:2606.04119v1 of this article is distributed under the Creative Commons Attribution 4.0 International licence, as recorded in the arXiv licence metadata for this exact version and shown on the abstract page https://arxiv.org/abs/2606.04119v1. A full rights notice is delivered at licenses/arxiv-2606.04119-CC-BY-4.0.txt. Characters: every retained excerpt is pure ASCII, so the delivered engine, XeLaTeX with its default Latin Modern text font, reports no missing character.
\begin{equation}  \label{eq:S-R-relation}
S_{M-1,L-1}(g;x,y)=R_{M,L}(g;x,y), \qquad M,L\in\mathbb{N }.
\end{equation}

\begin{proposition}
\label{prop:block} For every $N\in \mathbb{N}$, 
\begin{equation*}
\left\Vert F_{N}\right\Vert _{L^{2}(\mathbb{T}^{2})}=1.
\end{equation*}%
Moreover, let 
\begin{equation*}
\Omega _{0}:=\left[ \frac{1}{4},\frac{1}{2}\right] \times \left[ \frac{1}{8},%
\frac{1}{4}\right] \subset \mathbb{T}^{2}.
\end{equation*}%
There exist absolute constants $c_{\ast},c_0>0$ and $N_0\in\mathbb{N}$ such
that, for all $N\ge N_0$, 
\begin{equation}
\left|S_{\lfloor Ny\rfloor-1,\lfloor Nx\rfloor-1}(F_N;x,y)\right| \ge
c_{\ast}\log N, \qquad (x,y)\in\Omega_0,  \label{eq:block-large-S}
\end{equation}
and consequently 
\begin{equation*}
\mu_2\left\{(x,y)\in\mathbb{T}^2: \left|S_{\lfloor Ny\rfloor-1,\lfloor
Nx\rfloor-1}(F_N;x,y)\right| \ge c_{\ast}\log N\right\} \ge c_0.
\end{equation*}
One may take $c_{\ast}=1/(4\pi)$ and $c_0=1/32$.
\end{proposition}

\begin{proof}
The Fourier coefficients of $F_N$ are 
\begin{equation*}
\widehat F_N(m,n)= 
\begin{cases}
N^{-1}\mathrm{e}\!\left(-mn/N\right), & 0\le m,n\le N-1, \\ 
0, & \text{otherwise.}%
\end{cases}%
\end{equation*}
Thus Parseval's identity gives 
\begin{equation*}
\left\|F_N\right\|_{L^2(\mathbb{T}^2)}^2 =\sum_{m,n\in\mathbb{Z}%
}\left|\widehat F_N(m,n)\right|^2
=\sum_{m=0}^{N-1}\sum_{n=0}^{N-1}\frac1{N^2}=1.
\end{equation*}

Let 
\begin{equation*}
M:=\left\lfloor Ny\right\rfloor ,\qquad L:=\left\lfloor Nx\right\rfloor
,\qquad \alpha :=Nx-L\in \lbrack 0,1),\qquad \beta :=Ny-M\in \lbrack 0,1).
\end{equation*}%
Then, by the definition of the auxiliary sums $R_{M,L}$, 
\begin{equation}
R_{M,L}(F_{N};x,y)=\frac{1}{N}\sum_{a=0}^{M-1}\sum_{b=0}^{L-1}\mathrm{e}%
\!\left( ax+by-ab/N\right) .  \label{eq:R-first}
\end{equation}%
Summing first in $b$ gives 
\begin{equation*}
R_{M,L}(F_{N};x,y)=\frac{1}{N}\sum_{a=0}^{M-1}\mathrm{e}\!\left( ax\right) 
\frac{1-\mathrm{e}\!\left( L(y-a/N)\right) }{1-\mathrm{e}\!\left(
y-a/N\right) }.
\end{equation*}%
Put $a=M-r$, $r=1,\dots ,M$. Since 
\begin{equation*}
y-\frac{M-r}{N}=\frac{\beta +r}{N},\qquad M=Ny-\beta ,\qquad L=Nx-\alpha ,
\end{equation*}%
we obtain the exact identity 
\begin{equation}
R_{M,L}(F_{N};x,y)=\frac{\mathrm{e}\!\left( Nxy\right) }{N}\sum_{r=1}^{M}%
\frac{\mathrm{e}\!\left( -(\beta +r)x\right) -\mathrm{e}\!\left( -\alpha
(\beta +r)/N\right) }{1-\mathrm{e}\!\left( (\beta +r)/N\right) }.
\label{eq:exact-identity}
\end{equation}

Fix $(x,y)\in\Omega_0$. Then $N/8-1\le M\le N/4$. For $1\le r\le M$ set 
\begin{equation*}
\theta_r:=\frac{\beta+r}{N}.
\end{equation*}
For $N\ge8$ we have $0<\theta_r\le3/8$. We next justify the two uniform
estimates that will be used in \eqref{eq:exact-identity}.

First, consider the reciprocal of $1-\mathrm{e}\!\left(\theta\right)$. As $%
\theta\to0$, 
\begin{equation*}
\mathrm{e}\!\left(\theta\right)=e^{2\pi i\theta} =1+2\pi
i\theta-2\pi^2\theta^2+O(\theta^3),
\end{equation*}
and therefore 
\begin{equation*}
1-\mathrm{e}\!\left(\theta\right) =-2\pi i\theta+2\pi^2\theta^2+O(\theta^3)
=-2\pi i\theta\left(1+i\pi\theta+O(\theta^2)\right).
\end{equation*}
Thus 
\begin{equation*}
\begin{aligned} \frac1{1-\et{\theta}} &=\frac1{-2\pi i\theta}
\left(1+i\pi\theta+O(\theta^2)\right)^{-1} \\ &=\frac{i}{2\pi\theta}
\left(1-i\pi\theta+O(\theta^2)\right) \\
&=\frac{i}{2\pi\theta}+\frac12+O(\theta). \end{aligned}
\end{equation*}
Hence the function 
\begin{equation*}
G(\theta):=\frac1{1-\mathrm{e}\!\left(\theta\right)}-\frac{i}{2\pi\theta}
\end{equation*}
satisfies 
\begin{equation*}
G(\theta)=\frac12+O(\theta) \qquad (\theta\to0).
\end{equation*}
The singular terms in the two summands defining $G$ have therefore
cancelled. Defining 
\begin{equation*}
G(0):=\frac12
\end{equation*}
makes $G$ continuous at $0$. On $(0,3/8]$ the function is continuous as
well, since $\mathrm{e}\!\left(\theta\right)\neq1$ there. Consequently $G$
is continuous on the compact interval $[0,3/8]$, and so there exists an
absolute constant $C_G$ such that 
\begin{equation*}
|G(\theta)|\le C_G, \qquad 0\le \theta\le \frac38.
\end{equation*}
It follows that, uniformly for all $r$ with $1\le r\le M$, 
\begin{equation*}
\begin{aligned} \frac1{N(1-\et{\theta_r})}
&=\frac1N\left(\frac{i}{2\pi\theta_r}+G(\theta_r)\right) \\ &=\frac{i}{2\pi
N\theta_r}+O\!\left(\frac1N\right) \\
&=\frac{i}{2\pi(\beta+r)}+O\!\left(\frac1N\right), \end{aligned}
\end{equation*}
because $N\theta_r=\beta+r$. The constant in this $O(1/N)$ term is
independent of $r$, $N$, $\alpha$, and $\beta$.

Second, since $0\leq \alpha <1$, Taylor's formula for the exponential gives 
\begin{equation*}
\mathrm{e}\!\left( -\alpha \theta _{r}\right) =e^{-2\pi i\alpha \theta
_{r}}=1+O(\alpha \theta _{r})=1+O\!\left( \frac{\beta +r}{N}%
\right) ,
\end{equation*}%
again uniformly in $r$. To combine these estimates, write 
\begin{equation*}
A_{r}:=\mathrm{e}\!\left( -(\beta +r)x\right) -1,\qquad \delta _{r}:=\mathrm{%
e}\!\left( -\alpha \theta _{r}\right) -1.
\end{equation*}%
Then $|A_{r}|\leq 2$ and $\delta _{r}=O(\theta _{r})=O((\beta +r)/N)$.
Therefore 
\begin{equation*}
\begin{aligned} &\frac{\et{-(\beta+r)x}-\et{-\alpha(\beta+r)/N}}
{N(1-\et{(\beta+r)/N})} \\ &\quad = (A_r-\delta_r)
\left(\frac{i}{2\pi(\beta+r)}+O\!\left(\frac1N\right)\right) \\ &\quad
=\frac{i}{2\pi}\frac{\et{-(\beta+r)x}-1}{\beta+r} +O\!\left(\frac1N\right).
\end{aligned}
\end{equation*}
Thus 
\begin{equation*}
\frac{\mathrm{e}\!\left( -(\beta +r)x\right) -\mathrm{e}\!\left( -\alpha
(\beta +r)/N\right) }{N(1-\mathrm{e}\!\left( (\beta +r)/N\right) )}=\frac{i}{%
2\pi }\frac{\mathrm{e}\!\left( -(\beta +r)x\right) -1}{\beta +r}+O\!\left( 
\frac{1}{N}\right)
\end{equation*}%
uniformly for $1\leq r\leq M$. Summing over $r$ and using $M\leq N/4$, %
\eqref{eq:exact-identity} yields 
\begin{equation}
R_{M,L}(F_{N};x,y)=\frac{i\mathrm{e}\!\left( Nxy\right) }{2\pi }%
\sum_{r=1}^{M}\frac{\mathrm{e}\!\left( -(\beta +r)x\right) -1}{\beta +r}%
+O(1),  \label{eq:R-asymptotic}
\end{equation}%
where the $O(1)$ term is absolute and uniform on $\Omega _{0}$.
Consequently, 
\begin{equation}
\begin{aligned} \left|R_{M,L}(F_N;x,y)\right| &\ge
\frac1{2\pi}\sum_{r=1}^{M}\frac1{\beta+r} -\frac1{2\pi}\left|\sum_{r=1}^{M}
\frac{\et{-(\beta+r)x}}{\beta+r}\right|-C_1. \end{aligned}
\label{eq:R-lower-start}
\end{equation}

We claim that 
\begin{equation}  \label{eq:osc-bound}
\left|\sum_{r=1}^{M}\frac{\mathrm{e}\!\left(-(\beta+r)x\right)}{\beta+r}%
\right| \le C_2, \qquad x\in[1/4,1/2],\quad \beta\in[0,1),\quad M\in\mathbb{N%
}.
\end{equation}
Indeed, 
\begin{equation*}
\sum_{r=1}^{M}\frac{\mathrm{e}\!\left(-(\beta+r)x\right)}{\beta+r} =\mathrm{e%
}\!\left(-\beta x\right)\sum_{r=1}^{M}\frac{\mathrm{e}\!\left(-rx\right)}{%
\beta+r}.
\end{equation*}
Let 
\begin{equation*}
A_m(x):=\sum_{r=1}^{m}\mathrm{e}\!\left(-rx\right).
\end{equation*}
Since $x\in[1/4,1/2]$, 
\begin{equation*}
\left|A_m(x)\right| =\left|\mathrm{e}\!\left(-x\right)\frac{1-\mathrm{e}%
\!\left(-mx\right)}{1-\mathrm{e}\!\left(-x\right)}\right| \le \frac2{\left|1-%
\mathrm{e}\!\left(-x\right)\right|} \le \frac2{2\sin(\pi/4)}=\sqrt2.
\end{equation*}
Abel summation gives 
\begin{equation*}
\sum_{r=1}^{M}\frac{\mathrm{e}\!\left(-rx\right)}{\beta+r} =\frac{A_M(x)}{%
\beta+M} +\sum_{r=1}^{M-1} A_r(x)
\left(\frac1{\beta+r}-\frac1{\beta+r+1}\right),
\end{equation*}
and hence 
\begin{equation*}
\left|\sum_{r=1}^{M}\frac{\mathrm{e}\!\left(-rx\right)}{\beta+r}\right| \le 
\frac{\sqrt2}{\beta+M} +\sqrt2\sum_{r=1}^{M-1}
\left(\frac1{\beta+r}-\frac1{\beta+r+1}\right) \le 2\sqrt2.
\end{equation*}
This proves \eqref{eq:osc-bound}.

On the other hand, 
\begin{equation*}
\sum_{r=1}^{M}\frac{1}{\beta +r}\geq \log N-(\log 8+1).
\end{equation*}%
Combining this with \eqref{eq:R-lower-start} and \eqref{eq:osc-bound}, we
get 
\begin{equation*}
\left\vert R_{M,L}(F_{N};x,y)\right\vert \geq \frac{1}{2\pi }\log
N-C_{3},\qquad (x,y)\in \Omega _{0},
\end{equation*}%
with an absolute constant $C_{3}$. Choosing $N_0$ sufficiently large gives 
\begin{equation*}
\left|R_{M,L}(F_N;x,y)\right| \ge \frac{1}{4\pi}\log N, \qquad
(x,y)\in\Omega_0, \qquad N\ge N_0.
\end{equation*}
Since $F_N$ has spectrum in the nonnegative quadrant, the relation %
\eqref{eq:S-R-relation} gives 
\begin{equation*}
S_{\lfloor Ny\rfloor-1,\lfloor Nx\rfloor-1}(F_N;x,y) =R_{\lfloor
Ny\rfloor,\lfloor Nx\rfloor}(F_N;x,y).
\end{equation*}
This proves \eqref{eq:block-large-S}. The measure estimate follows because $%
\Omega_0$ has measure $1/32$. Thus one may take $c_0=1/32$.
\end{proof}

\end{document}
